\chapter{Taylor級数}
    \section{\texorpdfstring{$\NapierE^{-x^2}$}{e\string^(-x\string^2)}は全域でTaylor展開可能}
        \begin{proof}
            \quad\par
            $f(x) = \NapierE^{-x^2}\;(x\in\realNumbers)$とする。
            これ自体は$f$が複素関数として全平面で解析的であることを示して結果を実数に限定することで容易に示せるが、ここでは敢えて実解析の範囲で剰余項の評価を評価することでTaylor展開可能であることを示す。
            $f$の$n$階微分は$x$の$n$次多項式と$\NapierE^{-x^2}$の積になり、多項式の係数の最大の絶対値が$2^n$であることは容易に確かめられる。
            よって
            \[ |f^{(n)}(x)| \leq \NapierE^{-x^2}\sum_{i=0}^n 2^n |x|^i \]
            右辺は$x$を固定するごとに高々定数の$n$乗なので$n!$で除した結果は$n\to\infty$で0に収束する。
        \end{proof}
